\documentclass[11pt]{amsart}
\usepackage[T1]{fontenc}
\usepackage{amsmath,amssymb,mathtools,mathrsfs,booktabs,longtable,microtype}
\usepackage[margin=1.08in]{geometry}
\usepackage[hidelinks]{hyperref}
\hypersetup{pdftitle={General Theta Foundations I: Acquired Geometry and Causal Resource Transfer},pdfauthor={Qian Qi}}
\pdfinfoomitdate=1
\pdftrailerid{}
\pdfsuppressptexinfo=15
\newtheorem{theorem}{Theorem}[section]
\newtheorem{proposition}[theorem]{Proposition}
\newtheorem{lemma}[theorem]{Lemma}
\newtheorem{corollary}[theorem]{Corollary}
\theoremstyle{definition}
\newtheorem{definition}[theorem]{Definition}
\newtheorem{example}[theorem]{Example}
\theoremstyle{remark}
\newtheorem{remark}[theorem]{Remark}
\newcommand{\E}{\mathbb E}
\newcommand{\PP}{\mathbb P}
\newcommand{\R}{\mathbb R}
\newcommand{\N}{\mathbb N}
\newcommand{\TV}{\operatorname{TV}}
\newcommand{\law}{\operatorname{Law}}
\newcommand{\osc}{\operatorname{osc}}
\newcommand{\dist}{\operatorname{dist}}
\newcommand{\cp}{\mathrm{cp}}
\newcommand{\on}{\mathrm{on}}
\newcommand{\one}{\mathbf 1}
\newcommand{\norm}[1]{\left\lVert#1\right\rVert}
\newcommand{\abs}[1]{\left\lvert#1\right\rvert}
\title{General Theta Foundations I:\\Acquired Geometry and Causal Resource Transfer}
\author{Qian Qi}
\date{7 October 2026}
\subjclass[2020]{Primary 60G35, 62B05; Secondary 94A34, 93E11, 28A80}
\keywords{causal experiment, predictive quotient, acquired geometry, continuation information, adaptive acquisition, finite memory}
\begin{document}
\begin{abstract}
We study when the predictions acquired by a causal experiment can be retained with finite resources. A transfer theorem combines a constructive block-stability upper bound with lower bounds at every dominated continuation cut. The latter bounds concern the actual distribution of conditional future-response functions, rather than only the terminal predictive marginal. They give a necessary obstruction to matching checkpoint and online prediction. We also prove a matching acquisition--memory law, optimized over adaptive sensing policies, for anisotropic nonhomogeneous binary-cylinder experiments. This class includes smooth, singular and mixed product laws without requiring a limiting dimension. An independent erasure interface realizes a linear deficiency term in the same scored task, and a finite-bit implementation accounts separately for retained data, program, workspace and time. A second raw-kernel realization gives the sharp persistent-state order for the usual squared filtering loss in sparse Gaussian hidden-state models. Stopped comparison and resource-certified causal morphisms preserve one common task. The general stability certificate is sufficient; the continuation obstruction is necessary on its stated class of cuts, not a classification of all causal encoders.
\end{abstract}
\maketitle
\section{The experimental problem}\label{sec:problem}
A causal experiment does not present an observer with a parameter manifold. It presents preparations, admissible interventions and reports. Geometry becomes relevant only after specifying which future reports are to be predicted and which information the observer has actually acquired. There is a further obstruction: a small terminal codebook need not be recursively selectable after earlier reports have been discarded. The mother problem considered here is to connect these two facts to executable comparison of experiments with an explicit resource account.

The central result, Theorem~\ref{thm:transfer}, has two sides. A static cover preserving a budget-independent relation, together with physical-law block stability and robust candidate moments, constructs a finite-state predictor rounded after every report. In the other direction, a dominated continuation cut turns the conditional responses to future reports into a random element of a separable function space. Its actual acquired law supplies a lower bound on every retained alphabet crossing that cut. Terminal quantization is the last member of this family of bounds, not a replacement for it. The distinction produces an explicit online/checkpoint separation, using a classical random-access argument.

A separate difficulty concerns acquisition itself. A geometric lower bound for one fixed exploration does not determine the value of optimizing the exploration. Theorem~\ref{thm:adaptive} addresses this optimization for a class of progressive binary sensors. A single greedy scale sequence controls both the optimal acquisition allocation and the physical-law small-ball bound. Consequently adaptive sensing, online compression and terminal checkpoint compression have the same order under joint acquisition and retained-state constraints. Ratios may vary arbitrarily with scale and coordinate; some factors may be Lebesgue intervals while others are singular. Theorem~\ref{thm:joint} places acquisition, memory, calibration and an attained simulation defect in one common task. The second realization, Theorem~\ref{thm:brier}, concerns squared filtering loss for a primitive sparse hidden-state model with informative missingness and Gaussian reports.

These conclusions have different logical statuses. The block construction extends earlier sufficient results. Continuation quantization is a general obstruction with a product-domination hypothesis, not a sufficiency theorem for recursive realizability. The adaptive theorem optimizes the sensing class below, not arbitrary interventions. Quantization, conditional projection, random-access entropy and comparison inequalities are credited as antecedents. The technical companion~\cite{companion} retains the singular nonreset, stopped, calibration, output-grid and interface results with their original hypotheses.

\subsection{Kernels, interfaces and resources}
All measurable spaces are standard Borel. An experiment is specified by a preparation $\mu_\lambda$, hidden states $X_t$, actions $A_t$, reports $Y_{t+1}$, and jointly measurable probability kernels
\begin{equation}\label{eq:kernel}
 K_t^\lambda(dx',dy,dc\mid x,a).
\end{equation}
Here $c$ records nonnegative raw-call, time and other declared resource increments. Positivity means nonnegativity and normalization. Dirac operations, failure symbols and missing reports are allowed. A visible history $h_t$ consists of the preparation report and the intervening actions, reports and public costs. Legal action sets have a declared measurable graph. Strategies are parameter-independent Borel kernels on those sets; stopping rules use the visible filtration. They cannot consult future reports or the unknown $\lambda$.

A continuation interface is a measurable function $v_t(h_t)$ which specifies legal future actions, available resources and any publicly supplied phase. A predictive realization is a standard Borel state $S_t=s_t(h_t)$ determining that interface, every declared future-test expectation, and jointly Borel report and update kernels $J_t,F_t$. When a phase is kept separate, write the state space as a measurable disjoint union $S=\coprod_{v\in V}S_v$ over a \emph{finite} public phase set. All comparison pairs at a given time lie in the same fiber $S_v$. Updates map both members into the same next fiber, whose phase is a declared function of the old phase, action and common report. No unretained data-dependent phase is public for free. An unbounded or data-dependent external interface requires a separately stated conditional theorem and is not covered by the alphabet count below.

A predictor retains at most $M$ data labels at each declared cut. Its update uses only its current retained tuple, current action/report and allowed fresh randomness. Any controller, simulator or clock register consulted later is counted separately or included in the label tuple. Temporary workspace must be erased before the next cut; program constants must be known, parameter-independent data with a charged description. The notation $M$ is an alphabet cardinality, not total arithmetic space. If a public finite phase has $Q$ states, the full serial cardinality is at most $MQ$; a lower bound at a cut uses the \emph{entire} accessible data-dependent tuple. Deterministic time is fixed in the main prediction statements and is not a hidden observation channel.

For the general transfer theorem fix a preparation and a legal exploration independently of the prediction encoder. Algorithm coins are independent of that experiment. In the adaptive acquisition theorem the exploration is instead optimized together with the encoder; its lower bounds are proved anew and do not condition an encoder-dependent law as though it were fixed.

\subsection{Executable tests and a measurable predictive quotient}
An executable test is a legal finite continuation followed by a bounded function of its scored reports. Two histories are equivalent if their continuation interfaces and all test expectations agree. A set-theoretic quotient need not be a suitable measurable state. We record a sufficient realization result.

\begin{proposition}[Continuous-instrument realization]\label{prop:quotient}
Let $\mathcal B$ be compact metrizable and let $\mathcal V\subset C(\mathcal B)$ be the separable closed linear span of constants and the continuous evaluations of a determining family of executable tests and interface coordinates. All kernel, density, update and instrument versions below are jointly Borel also in the action variable on the declared standard Borel action--report domain. For each fixed action, let its report space be Polish, with reference measure $m_a$, and let $g_a(\pi,y)$ be a jointly continuous density. On the open set $D_a=\{(\pi,y):g_a(\pi,y)>0\}$, with the relative product topology, suppose the normalized update $B_a$ is continuous. For every $f\in\mathcal V$, require the instrument evaluation $g_a(\pi,y)f(B_a(\pi,y))$, with its specified value at density zero, to belong to $\mathcal V$ as a function of $\pi$ and to be jointly continuous. Then the countable evaluation image is a compact metrizable predictive quotient. Report laws and updates descend, and the latter are continuous on the descended open positive-density domain. Every standard Borel statistic measurably determining all these evaluations measurably determines the quotient.
\end{proposition}
\begin{proof}
Normalize a countable dense family of evaluations and map $\mathcal B$ into $[-1,1]^{\N}$. The image $S$ is compact, hence Borel. Equality of its coordinates gives equality on $\mathcal V$ by uniform approximation and hence on the determining tests. Conversely, equality of all future tests and interface coordinates gives equality on their closed linear span. Thus the fibers are exactly predictive equivalence classes, not a finer partition carrying unrelated coordinates. Taking $f=1$ shows that density values agree on each fiber. Dividing equal instrument values by that positive common density shows that the evaluation of the updated belief agrees on the fiber.

The evaluation map is a continuous surjection from a compact space to a Hausdorff space, hence is perfect. Its product with the report identity is also a quotient map: closedness follows by taking a convergent subnet in the compact belief domain whenever an image net converges. The instrument coordinates and their positive-density ratios therefore descend continuously. The positive-density domain is open in $S\times Y$; no update is defined by division at a null report. For measurability jointly across actions, the evaluation map admits a Borel section. Indeed choose nested finite closed covers of the compact domain with diameters tending to zero and, at each stage, the first member meeting the fiber and the previously selected members. Each such test is Borel because the image of the resulting fixed compact intersection is closed. The selected nested intersections have a unique limit in the fiber, and their shrinking diameters make that limit a Borel function of the image point. Composing the assumed joint versions with this section supplies joint Borel versions of all descended kernels. A common jointly Borel null-report extension, when needed on wrong candidates, is additional data.

For minimality, collect the measurable coordinate factors through the given statistic into a measurable product map. It takes values in the Borel set $S$ on the domain in question; extend outside that set by a fixed point. Countably many almost-sure factors can be imposed on a single common full-measure history set. This proves measurable, rather than merely set-theoretic, minimality.
\end{proof}

Applications may verify a predictive realization directly. The metric used for a quantitative estimate need only induce its standard Borel structure; it need not be the compact topology in Proposition~\ref{prop:quotient}.

\subsection{One prediction task}
Let $Y$ be a bounded terminal vector in a finite-dimensional Euclidean space with a specified weighted norm. A finite independent terminal probe index is interpreted as the vector of its separately executable conditional means, with the probe probabilities as weights; it is not a simultaneous counterfactual measurement. At the decision cut $T$ put
\begin{equation}\label{eq:task}
 P_T=\E[Y\mid\mathcal H_T],\qquad B_T=\E\norm{Y-P_T}^2,
 \qquad \nu_T=\law(P_T).
\end{equation}
For a finite Borel measure $\nu$ on a separable Hilbert space define
\begin{equation}\label{eq:quantization}
 q_M(\nu)=\inf_{c_1,\ldots,c_M}\int\min_i\norm{x-c_i}^2\,\nu(dx).
\end{equation}
The measure may be atomic, singular or a subprobability. A checkpoint algorithm inspects $\mathcal H_T$ once before retaining $M$ labels. An online algorithm retains at most $M$ labels at every earlier cut and cannot reread a report. In both cases the terminal outcome is revealed only after the prediction. Denote the optimal risks by $R_\cp(T,M)$ and $R_\on(T,M)$.

\section{Acquired continuation geometry and causal transfer}\label{sec:transfer}
We first specify the lower certificate, which depends only on the physical experiment and its fixed continuation, not on a proposed encoder.

\begin{definition}[Dominated continuation cut]\label{def:cut}
At a deterministic cut $s\le T$, let $H$ be the actual prefix and $U$ the complete physical pre-decision suffix. The complete decision history is generated by $(H,U)$. The fixed exploration supplies their joint law $\Lambda(dh,du)$. A dominated continuation certificate consists of a probability measure $\eta$ on the suffix space and a finite measure $\underline\mu$ on prefix histories such that
\begin{equation}\label{eq:domination}
 \underline\mu(dh)\eta(du)\le\Lambda(dh,du).
\end{equation}
Choose a bounded jointly Borel version $z(h,u)=\E[Y\mid H=h,U=u]$ and set
\[
 Z_s(h)=z(h,\cdot)\in L^2(\eta;\R^q),\qquad
 \xi_s=(Z_s)_*\underline\mu.
\]
The task norm supplies the inner product in this function space. All post-cut information supplied to a decoder is contained in $U$ and independent algorithm coins; no earlier tape or additional data-dependent public register remains outside its counted label.
\end{definition}

The product inequality is a raw-law minorization, not independence of the actual suffix. For example a density $j(u\mid h)\ge c$ relative to $\eta$ on an actual prefix event $E$ gives $\underline\mu=c\,\law(H)|_E$. An independent suffix has $c=1$. In the absence of such a nonzero certificate the corresponding bound is zero, not an assertion that the experiment is easy. Domination ensures that versions of $z$ give the same element for $\underline\mu$-almost every $h$. Standard Borel spaces have countably generated sigma fields modulo a fixed probability, so $L^2(\eta)$ is separable. Approximation of a bounded jointly measurable function by simple rectangle functions, followed by dominated convergence, makes $h\mapsto Z_s(h)$ a Borel map into this Hilbert space.

\begin{lemma}[Projection and the continuation cut]\label{lem:cut}
For the fixed task in \eqref{eq:task},
\begin{equation}\label{eq:checkpoint}
 R_\cp(T,M)=B_T+q_M(\nu_T).
\end{equation}
Every $M$-label online predictor and every certificate in Definition~\ref{def:cut} satisfy
\begin{equation}\label{eq:cutlower}
 R_\on(T,M)-B_T\ge q_M(\xi_s).
\end{equation}
Both assertions allow independent shared and private randomization. They use the same terminal task and baseline.
\end{lemma}
\begin{proof}
Conditional projection gives $\E\norm{Y-A}^2=B_T+\E\norm{P_T-A}^2$. Fix all algorithm coins, which are independent of the fixed exploration. A checkpoint readout has at most $M$ centers. Randomized readouts may alternatively be replaced by conditional means, and nearest-center selection is Borel. Taking approximating center lists proves \eqref{eq:checkpoint}; averaging independent seeds does not change $\nu_T$.

At cut $s$ write its retained label as $i(H)$. Give the suffix decoder all of $U$, even when its original online implementation would retain less. Conditional on fixed coins it has an output $a_{i(H)}(U)$ for at most $M$ Borel response functions. Projection and \eqref{eq:domination} give
\[
 R-B_T\ge\int\norm{z(h,u)-a_{i(h)}(u)}^2
                   \underline\mu(dh)\eta(du)
 \ge\int\min_i\norm{Z_s(h)-a_i}_{L^2(\eta)}^2\underline\mu(dh).
\]
Outputs can be projected into a fixed bounded convex set containing the conditional means, so all center functions belong to $L^2(\eta)$. The last expression is at least $q_M(\xi_s)$. Averaging the coins proves the assertion. No different continuation task is inserted into the loss, and no product structure of $\Lambda$ beyond the lower certificate is used.
\end{proof}

We next state a constructive upper certificate. It is convenient to allow a finite declared phase as in Section~\ref{sec:problem}; otherwise take one fiber. All constants below are uniform over phases and the specified strategies. A code has at most $M$ representatives in each phase. A cut lower must replace $M$ by the size of the full accessible tuple when its phase itself carries data.

Let $(S_v,d_v)$ be separable metric predictive fibers, with Borel envelope $w\ge1$. The update along a common report and physical action is $F_t$. A Borel relation $\mathcal A_v$ satisfies $x\in\mathcal A_v(x)$ and $w(z)\le w(x)$ for $z\in\mathcal A_v(x)$. It is fixed independently of the budget. Finite Borel covers have representatives $c_i^v$ and first-cell selectors $Q_M^v$ such that
\begin{equation}\label{eq:cover}
 Q_M^v(x)\in\mathcal A_v(x),\qquad
 d_v(x,Q_M^v(x))\le r_M w(x).
\end{equation}
For $s=km$, write $\mathscr F_{s,j}(z)$ for $j$ exact updates along the next physical word and the actions selected from that physical transcript. For every possible block-start past, true state $x$ and candidate $z$ in its named fiber, require one-step pathwise Lipschitz gain at most $L\ge1$ and
\begin{equation}\label{eq:pair}
 \E_x[d(\mathscr F_{s,m}(x),\mathscr F_{s,m}(z))^2\mid\mathcal H_s]
 \le\kappa d(x,z)^2,\qquad 0\le\kappa<1.
\end{equation}
The notation $d$ denotes the appropriate common fiber metric. For \emph{every} adapted sequence $z_0=z$, $v_j=F_{s+j}(z_{j-1},a_{s+j},Y_{s+j})$, $z_j\in\mathcal A(v_j)$, require
\begin{align}
 \E_x[w(z_m)^2\mid\mathcal H_s]&\le a w(z)^2+b,\label{eq:drift}\\
 \E_x[w(v_j)^2\mid\mathcal H_s]&\le A w(z)^2+B\quad(1\le j\le m),\label{eq:within}
\end{align}
where $0\le a<1$ and $b,A,B\ge0$. The conditional inequalities continue to hold when past independent algorithm coins are included. They quantify over arbitrary relation-preserving adapted perturbations, rather than assuming the performance of a preconstructed quantizer. Numerical implementations must still choose the same codebook and relation, with local error at most $r_Mw(v_j)+u$.

\begin{theorem}[Causal acquired-geometry and continuation transfer]\label{thm:transfer}
Fix a prepared experiment, a deterministic horizon and encoder-independent exploration. Suppose \eqref{eq:cover}--\eqref{eq:within} and a legal seed with $\E w(\widehat S_0)^2\le H_0$ and $\norm{d(S_0,\widehat S_0)}_2\le e_0$. Assume $P_T=f_T(S_T)$ with Lipschitz constant $L_f$ in the task norm. Put
\begin{align*}
 H&=\max\{H_0,b/(1-a)\},& H_*&=AH+B,\\
 \Lambda_j&=\sum_{i=0}^{j-1}L^i,& \Delta_M&=r_M\sqrt{H_*}+u,\\
 E_*&=\max\{e_0,\Lambda_m\Delta_M/(1-\sqrt\kappa)\},&
 E_M&=L^{m-1}E_*+\Lambda_{m-1}\Delta_M.
\end{align*}
For a one-fiber $M$-label model, or a model in which $M$ below counts the whole accessible tuple, define
\[
 W_M=\max\{q_M(\nu_T),\ \sup_{s,\,\text{certificates}}q_M(\xi_s)\}.
\]
There is a causal predictor quantized after every raw call, with terminal readout error at most $v$ in root mean square, such that
\begin{equation}\label{eq:central}
 \begin{gathered}
 R_\cp(T,M)=B_T+q_M(\nu_T),\\
 B_T+W_M\le R_\on(T,M)
       \le B_T+(L_fE_M+v)^2.
 \end{gathered}
\end{equation}
For phase-indexed covers the upper uses the stated $M$ labels per phase and their charged full tuple; the lower always uses that tuple's cardinality.

For any actual measure $\xi$, either $\nu_T$ or a dominated continuation measure, a submeasure of mass $\zeta>0$ with
\begin{equation}\label{eq:smallball}
 \sup_c\underline\xi(B(c,s_M))\le\zeta/(2M)
\end{equation}
contributes at least $\zeta s_M^2/2$ to the lower. Thus a matching actual terminal witness with $s_M\asymp r_M$, mass bounded below, $e_0,u,v=O(r_M)$ and fixed certificate constants gives matching checkpoint and online excess of order $r_M^2$. A larger continuation width instead certifies a separation from the terminal checkpoint curve. All terms refer to the one fixed task and actual preparation.
\end{theorem}
\begin{proof}
The exact selector in \eqref{eq:cover} is Borel and yields the recursion $\widehat S_t=Q_M(F_t(\widehat S_{t-1},a_t,Y_t))$. The numerical routine obeys the same relation. Equation~\eqref{eq:drift} proves inductively that the envelope second moment is at most $H$ at block endpoints; \eqref{eq:within} bounds each pre-rounding moment by $H_*$. Hence each local insertion of error has $L^2$ norm at most $\Delta_M$.

Start an exact shadow at $\widehat S_s$ for a block and feed it the same physical word. Telescoping the local insertions, with all subsequent pathwise gains included, gives
\[
 d(\mathscr F_{s,m}(\widehat S_s),\widehat S_{s+m})
 \le\sum_{j=1}^m L^{m-j}d(v_j,z_j).
\]
Conditional use of \eqref{eq:pair}, then Minkowski, shows that the root errors satisfy
\[
 e_{s+m}\le\sqrt\kappa e_s+\Lambda_m\Delta_M,
 \qquad e_{s+j}\le L^je_s+\Lambda_j\Delta_M\quad(0\le j<m).
\]
The first preserves $e_{km}\le E_*$ and the second gives $e_T\le E_M$. The shadow is a proof device: no block of reports, real posterior or second state is retained by the machine. Lipschitz readout, root-error addition and conditional projection prove the upper.

Lemma~\ref{lem:cut} gives all the lower bounds simultaneously for this same experiment. Their supremum remains a lower bound. Finally $M$ open balls satisfying \eqref{eq:smallball} cover at most half the submeasure. Outside their union the squared distance from every list is at least $s_M^2$, proving the stated contribution. This is valid in the continuation Hilbert space as well as in a finite-dimensional score space. The matching assertions follow directly; no ambient-dimensional volume has been assumed.
\end{proof}

\begin{corollary}[A necessary compatibility condition]\label{cor:necessary}
For a sequence of fixed-task experiments and budgets, an asserted bound
$R_\on(T,CM)-B_T\le C' q_M(\nu_T)$ implies
$W_{CM}\le C' q_M(\nu_T)$, with the same specified full-state cut. A family for which this inequality fails cannot have that claimed online/checkpoint matching, even if its terminal marginal has small covering number.
\end{corollary}
\begin{proof}
Apply the lower side of \eqref{eq:central}, or Lemma~\ref{lem:cut} alone when no stability certificate is available.
\end{proof}

\begin{example}[A quantitative delayed-query separation]\label{ex:query}
Prepare $d$ independent fair bits $X_1,\ldots,X_d$ and report them in $d$ paid calls. After the last retained cut report an independent uniform index $J\in\{1,\ldots,d\}$. The task is to predict $Y=X_J$ before its terminal audit. A checkpoint reading the full history after $J$ needs two labels for zero risk. If an online predictor had at most $M$ labels immediately before $J$, then
\begin{equation}\label{eq:randomaccess}
 R_\on\ge\tfrac12 h_2^{-1}\bigl((1-\log_2 M/d)_+\bigr),
\end{equation}
where $h_2^{-1}$ is the inverse binary entropy on $[0,1/2]$. In particular retaining at most a fixed fraction of $d$ bits gives a positive error bounded away from zero. Exact prediction requires $M\ge2^d$ and storing the bit string attains it.
\end{example}
\begin{proof}
The continuation measure is the law of $j\mapsto X_j$ in the uniform $L^2$ norm. For an independent shared seed $R$ and retained label $C$, $I(X;C\mid R)\le\log_2 M$. Independence of the bits and conditional subadditivity give
$\sum_j H(X_j\mid C,R)\ge d-\log_2M$.
Let $e$ be the average posterior bit-classification error. Concavity of $h_2$ gives $h_2(e)\ge(1-\log_2M/d)_+$. A Bernoulli posterior with smaller probability $p\le1/2$ has variance $p(1-p)\ge p/2$. Conditional projection proves \eqref{eq:randomaccess}. Zero error forces distinct positive-probability labels for all strings. This is the classical entropy mechanism behind random-access bounds~\cite{nayak}, used here to distinguish two information cuts of a single prediction task, not claimed as a new coding inequality.
\end{proof}

\begin{proposition}[Flagged violations of an admissibility relation]\label{prop:repair}
Suppose a finite implementation agrees with a routine satisfying the upper hypotheses until a flagged first failure, and, before that failure, its conditional probability of failing at step $t$ is at most $\beta_t$. A repair routine using the same codebook can always be selected on a failure to satisfy the original relation and local-error bound. For a loss in $[0,L_*]$, the actual implementation has risk at most the repaired upper plus $L_*\min(1,\sum_{t\le T}\beta_t)$.
\end{proposition}
\begin{proof}
Drive the repaired routine by the same physical reports and coins; at a failure choose the certified selector instead. Encoder-independent exploration makes the physical law unchanged. The two retained trajectories agree until the first failure. A first-failure union bound gives its probability at most $\sum_t\beta_t$, and bounded loss gives the assertion. The reference estimate is unconditional; it is not incorrectly conditioned on the absence of failure. The proposition weakens implementation accuracy to a probabilistic certificate, but does not assert the existence of global relation-preserving covers from local data alone.
\end{proof}

\section{Adaptive acquisition on nonhomogeneous cylinders}\label{sec:adaptive}
The cut argument concerns a fixed exploration. We now optimize acquisition itself. The raw interface in this section is progressive sensing: a call selects a coordinate and reveals its next unread binary digit. Neither a coordinate value nor an infinite tail is supplied as a pre-decision report. This interface is appropriate to a multiresolution measurement, and its restrictions are part of the theorem.

Fix $d\ge1$, positive weights $\omega_j$ with $\sum_j\omega_j=1$, and $0<a\le1/2$. Let
\begin{equation}\label{eq:scales}
 a\le r_{j,l}\le1/2,\qquad
 \ell_{j,0}=1,\quad \ell_{j,k}=\prod_{l=1}^k r_{j,l},\quad
 b_{j,l}=(1-r_{j,l})\ell_{j,l-1}.
\end{equation}
Prepare independent fair bits $\zeta_{j,l}$ and physical coordinates
\begin{equation}\label{eq:cylinders}
 X_j=\sum_{l\ge1}b_{j,l}\zeta_{j,l},\qquad
 \norm{x}_\omega^2=\sum_{j=1}^d\omega_jx_j^2.
\end{equation}
The series converges absolutely and $\sum_l b_{j,l}=1$. After preparation a legal action is either $\mathrm{read}(j)$, which returns the next unread bit of coordinate $j$, or $\mathrm{idle}$, which returns a fixed symbol. Each costs one raw call and one unit of physical acquisition time. The apparatus advances the chosen coordinate's internal read position. That position is part of the physical state, not a tape of past outcomes available to the observer. A policy may choose its actions from its permitted retained history and randomness. It cannot query a digit twice, request the hidden tail in one call, or obtain the terminal audit before predicting.

There are $n$ pre-decision opportunities; shorter acquisition rules are padded with paid idles to this deterministic horizon. A terminal audit returns $X+s\delta\mathbf 1$, with unknown $s\in\{-1,1\}$ and $0\le\delta\le1/8$, only \emph{after} the vector forecast. This is the standard weighted squared reconstruction task. The terminal audit costs one call and does not furnish an input to the prediction algorithm. Preparation also costs one call, so $N=n+2$. Actions and stopping cannot reveal $s$. Decisions may be restricted to the compact box $[-1/8,9/8]^d$, without increasing risk. The oracle risk is zero and the loss is bounded by $25/16$.

A complete history knows a finite prefix in each coordinate. Write $k_j(h)$ for its number of reads and $u_j(h)$ for that prefix. Its conditional mean is
\begin{equation}\label{eq:actualcenter}
 P_j(h)=\tfrac12+\sum_{l\le k_j(h)}b_{j,l}(u_{j,l}(h)-\tfrac12).
\end{equation}
The prefixes and read positions determine every legal future test by direct kernel integration. Conversely, the continuation which reads and scores a chosen coordinate determines its physical conditional law and read position; including the interface gives a predictive realization. For an adaptive policy the nominal acquired score law is exactly
\begin{equation}\label{eq:adaptiveactual}
 \nu_\alpha=\sum_{h}\PP_\alpha(H=h)\delta_{P(h)}.
\end{equation}
The sum is over the finitely many action/report histories of the fixed horizon. Its weights are the products of the policy probabilities and the factors $1/2$ for the actually read digits. Randomizing the policy averages these weights. Prefixes are \emph{not} asserted to be uniform conditional only on the random depths: data-dependent choices may bias that conditioning.

\subsection{A common scale for acquisition and quantization}
Put $s_j(k)=\sqrt{\omega_j}\ell_{j,k}$. Starting from $k^0=0$, increment, at each step, a coordinate with largest $s_j(k_j)$, breaking ties by the smallest index. Set
\begin{equation}\label{eq:greedy}
 R_b=\max_j s_j(k_j^b),\qquad b=0,1,\ldots.
\end{equation}
The sequence is known from the experiment, not estimated from the observed bits.

\begin{lemma}[Greedy scale and product small balls]\label{lem:greedy}
For every integer $b\ge0$,
\begin{equation}\label{eq:minmaxscale}
 R_b=\min_{\substack{k_j\ge0\ \sum_jk_j\le b}}\max_j s_j(k_j),
 \qquad aR_b\le R_{b+1}\le R_b.
\end{equation}
If $\mu=\law(X)$ and $b_M=\lfloor\log_2M\rfloor$, then
\begin{equation}\label{eq:physicalquant}
 c_{a,d}R_{b_M}^2\le q_M(\mu)\le(d/4)R_{b_M}^2.
\end{equation}
The constants are independent of the weights, the scale sequence and $M$. In particular no limiting dimension is assumed.
\end{lemma}
\begin{proof}
Each list $s_j(0),s_j(1),\ldots$ decreases by a factor in $[a,1/2]$. An increment removes its currently largest remaining entry. The greedy rule therefore removes the $b$ largest entries in the union of the $d$ lists, with a fixed order for ties. Every allocation of $b$ increments removes a prefix of each list. Its maximum remaining entry is at least the $(b+1)$st entry of the merged order; greedy attains it. This proves the first equality, including unused increments by monotonicity. The largest entry removed at the next step is $R_b$, and its successor is at least $aR_b$; this gives the second assertion.

At the greedy allocation $k^b$, every coordinate with $k_j^b>0$ has final weighted cylinder length at least $aR_b$. Indeed its last parent was a maximum at an earlier greedy step, hence was at least the eventual maximum $R_b$. In a fixed coordinate all level-$k_j^b$ cylinders have that same length. They have disjoint interiors because each ratio is at most $1/2$; their endpoints have zero fair-product mass. A ball for $\norm{\cdot}_\omega$ of radius $aR_b/4$ has projection length at most $aR_b/2$ in each weighted coordinate. It meets at most three of these intervals in each active coordinate. Each product cylinder has mass $2^{-b}$. Consequently
\begin{equation}\label{eq:productsmallball}
 \sup_x\mu(B(x,aR_b/4))\le3^d2^{-b}.
\end{equation}
Coordinates never split contribute just one interval and do not weaken the bound.

Choose $L_d=\lceil\log_2(4\cdot3^d)\rceil$ and $b=b_M+L_d$. Since $M<2^{b_M+1}$, the union of $M$ such balls has mass at most $1/2$. Its complement contributes at least $a^2R_b^2/32$. Iterating \eqref{eq:minmaxscale} gives $R_b\ge a^{L_d}R_{b_M}$, proving the lower with $c_{a,d}=a^{2L_d+2}/32$.

For the upper use the conditional means of the $2^{b_M}$ product cylinders at allocation $k^{b_M}$. Within coordinate $j$ the conditional distribution lies in an interval of length $\ell_{j,k_j^{b_M}}$, so its variance is at most one quarter of that squared length. Summing the weighted variances gives at most $(d/4)R_{b_M}^2$. The centers are physical conditional means; no smooth reference density or Fisher matrix has been introduced.
\end{proof}

\begin{lemma}[Adaptive acquisition variance]\label{lem:acquisition}
For every causal adaptive sensing policy with at most $n$ reads, conditional on a complete realized history the unread digits remain independent and fair. In particular
\begin{equation}\label{eq:tailvariance}
 \frac{\ell_{j,k}^2}{16}\le
 V_j(k):=\frac14\sum_{l>k}b_{j,l}^2
 \le\frac{\ell_{j,k}^2}{4},
\end{equation}
and the full-history acquisition uncertainty satisfies
\begin{equation}\label{eq:adaptivelower}
 \E\norm{X-P(H)}_\omega^2
   =\E\sum_j\omega_j V_j(k_j(H))\ge R_n^2/16.
\end{equation}
This remains true if the sensing controller has unrestricted access to its past reports.
\end{lemma}
\begin{proof}
Condition also on independent policy coins. Initially the claim is the product preparation. A selected coordinate is a function of revealed data. Its next unread bit is fair and independent of all other unread bits. Revealing it, or performing an idle, preserves the induction. Conditional means and variances are therefore given by \eqref{eq:actualcenter} and the displayed tail series. Its first term has coefficient $b_{j,k+1}\ge\ell_{j,k}/2$, giving the lower in \eqref{eq:tailvariance}. The whole tail has range of length $\ell_{j,k}$, giving the upper. Each realized allocation has sum at most $n$, so \eqref{eq:minmaxscale} bounds its largest weighted length below by $R_n$. Sum the nonnegative variances and then average. Removing the conditioning on coins proves the claim for randomized policies as well.
\end{proof}

\begin{theorem}[Optimized acquisition--memory law]\label{thm:adaptive}
Let $N=n+2\ge2$, $M\ge1$, and $K=\min\{n,\lfloor\log_2M\rfloor\}$. Optimize over all parameter-independent adaptive policies for the raw interface above and all online predictors whose entire data-dependent retained tuple has at most $M$ values at every pre-decision cut. Define the analogous checkpoint optimum by allowing the prediction encoder one final access to the complete acquisition history. Then
\begin{equation}\label{eq:adaptivelaw}
 c_{a,d}\{R_K^2+\delta^2\}
 \le\mathcal R_\cp(N,M,\delta)
 \le\mathcal R_\on(N,M,\delta)
 \le C_{a,d}\{R_K^2+\delta^2\}.
\end{equation}
The risk is minimax in the inaccessible sign $s$. A nonadaptive greedy allocation, followed by paid idles, attains the upper while storing only its $K$ observed bits. All constants are uniform in horizon, weights and ratios at fixed $a,d$. The deterministic phase/program is separately charged; it cannot contain observed data.
\end{theorem}
\begin{proof}
Fix any admissible procedure and condition on its independent seed. Averaging the risks at the two signs gives
\[
 \tfrac12\sum_{s=\pm1}\E\norm{X+s\delta\mathbf1-A}_\omega^2
 =\E\norm{X-A}_\omega^2+\delta^2.
\]
Full-history projection and Lemma~\ref{lem:acquisition} make the first term at least $R_n^2/16$. Independently, at the terminal decision it has only $M$ possible vector centers. The law of the physical $X$ is unchanged by any sensing policy or independent seed. Its distortion is therefore at least $q_M(\mu)\ge c_{a,d}R_{b_M}^2$ by Lemma~\ref{lem:greedy}. This is an oracle lower witness, not a substitution of $\mu$ for the actual atomic acquired law \eqref{eq:adaptiveactual}. Since $R_b$ decreases, the larger of $R_n^2$ and $R_{b_M}^2$ is $R_K^2$. These lower bounds hold even for a full-history sensing controller and hence for both classes in the theorem.

For the upper read the $K$ coordinates of the greedy schedule in order. Retain the resulting binary word. At the $t$th read there are $2^t\le2^K\le M$ possible words; all remaining opportunities are paid idles and copy that word exactly. The deterministic phase determines which coordinate each bit belongs to. Output its conditional mean \eqref{eq:actualcenter}. Its nominal squared error is at most $(d/4)R_K^2$. Its unconditional mean error is zero, so each sign adds exactly $\delta^2$. Preparation, all reads and idles, and the terminal audit give exactly $N$ calls. This construction also proves that the upper does not assume a supplied recursively accurate real-valued state.
\end{proof}

For any fixed policy the exact checkpoint identity still separates acquisition and memory:
\begin{equation}\label{eq:policyspecific}
 R_{\cp,\alpha}(M,0)
 =\E\norm{X-P(H)}_\omega^2+q_M(\nu_\alpha).
\end{equation}
In proving \eqref{eq:adaptivelaw} we did not commute an infimum over policies with either summand. Instead, the variance lower holds at every realized allocation and the physical-law converse holds for every encoder.

If all ratios equal $r$ and the weights are $1/d$, greedy depths differ by at most one and $R_K^2$ is comparable, at fixed $d,r$, to $r^{2K/d}/d$. With $r=1/2$ this includes Lebesgue product laws; with $r<1/2$ it includes singular products. Mixed and oscillating sequences are covered by the exact scale sequence even when no single power-law dimension exists. This is a nonuniform geometry statement, not a claim of uniformity as $a\downarrow0$ or $d\to\infty$.

\subsection{A common task attaining the deficiency term}
An upper allowance for deficiency cannot force a positive error floor. We give a source--target pair which attains it, while preserving the optimized geometric part of the task.

Add an unknown fixed bit $z\in\{0,1\}$ to the preparation. Exactly one designated paid call, immediately after preparation, reports $z$ with probability $1-2\varepsilon$ and an erasure otherwise, where $0\le\varepsilon\le1/2$. No other pre-decision call depends on $z$. Then allow the same $n$ progressive opportunities. The terminal audit, only after the prediction, scores
\begin{equation}\label{eq:jointtask}
 \ell((X+s\delta\mathbf1,z),(A,c))
 =\tfrac12\norm{X+s\delta\mathbf1-A}_\omega^2+\tfrac12(z-c)^2.
\end{equation}
There are exactly $N=n+3$ raw calls. The exact-bit target changes only the designated report to an exact bit and has the same costs and scored outcome. An admitted simulator must complete this virtual bit report before processing the progressive commands and may not query the terminal audit early.

\begin{theorem}[A matched common acquisition--memory--defect profile]\label{thm:joint}
For $N\ge3$, $M\ge3$ and $K=\min\{N-3,\lfloor\log_2M\rfloor\}$, the source experiment just defined satisfies
\begin{equation}\label{eq:jointlaw}
 \mathcal R_\cp\asymp_{a,d}\mathcal R_\on
 \asymp_{a,d}R_K^2+\delta^2+\varepsilon,
\end{equation}
where the infima optimize the permitted adaptive sensing and finite-label prediction, and the supremum is over $(s,z)$. Its parameter-uniform causal deficiency to the exact-bit target, with the stated completion interface and common task, is exactly $\varepsilon$. The target-to-source direction is exact by independent erasure. These claims concern this specified family, not every experiment admitting a defect upper bound $\varepsilon$.
\end{theorem}
\begin{proof}
Put the uniform prior on $(s,z)$ for a lower bound on the supremum. Conditional on an erased bit report, all later pre-decision information is independent of $z$, so its Bayes squared error is $1/4$. This contributes $\varepsilon/4$ to the weighted task. The geometric component has the acquisition and physical-law memory lower bounds used in Theorem~\ref{thm:adaptive}; knowing an independent bit or erasure indicator cannot invalidate them. The sign average adds $\delta^2/2$. Taking the larger of the acquisition and memory terms proves the lower.

For the upper retain one of $0,1,\bot$ and $K'=\min\{n,\lfloor\log_2(M/3)\rfloor\}$ greedily acquired digits. The tuple has at most $3\cdot2^{K'}\le M$ values. Predict the observed bit, or $1/2$ on erasure, and use the conditional mean for $X$. The bit contribution is $\varepsilon/4$ for either $z$. Since $0\le K-K'\le2$ and each scale step changes by a factor at least $a$, $R_{K'}^2\le a^{-4}R_K^2$. This proves the upper uniformly in the nuisance parameters and in $N,M$.

To simulate the exact-bit target, keep the observed bit and replace an erasure by an independent fair bit. Its chance of a wrong virtual report is $\varepsilon$ for each $z$. Couple the physical coordinates, all later policy coins and reports until that discrepancy; on its complement the complete common-task experiments coincide. This proves complete-law defect at most $\varepsilon$, uniformly over the admitted target policies. Conversely the two source distributions at the designated completion cut have total variation $1-2\varepsilon$. Their minimax bit-testing error is $\varepsilon$, since the only common report is the erasure. Any simulated exact-bit report would be such a test. Thus its worst-parameter total-variation defect from the exact report is at least $\varepsilon$. The compulsory completion order prevents use of the later audit. Independent erasure of an exact bit implements the reverse direction with zero defect.
\end{proof}

\subsection{A closed finite-bit realization}
The preceding theorems are exact-label statements. Their physical coefficient sequences may be arbitrary real data; such data do not automatically have finite programs. The following result states a finite implementation on a precisely specified region instead of hiding real constants in the machine.

\begin{theorem}[Finite data, scratch, program and time]\label{thm:digital}
For the upper construction in Theorem~\ref{thm:adaptive}, fix its finite greedy schedule of length $K$ and suppose the read-only calibration description supplies dyadic approximations $\widehat b_{j,l}$ for the selected coefficients, with error at most $2^{-p}$, $p\ge1$. The schedule is supplied as certified finite data; no decidability of exact comparison of arbitrary computable reals is assumed. There is a deterministic finite-bit implementation with pre-decision report alphabet $\{0,1,\mathrm{idle}\}$, $K$ retained data bits, and
\begin{align}
 \Gamma&=(K+1)2^{-p},\label{eq:digitalerror}\\
 W&\le C\{p+\log_2(K+d+2)\},\label{eq:workspace}\\
 L_{\rm prog}&\le C(K+d)\{p+\log_2(K+d+N+2)\},\label{eq:program}\\
 T_{\rm bit}&\le C\{N+d(K+d)(p+\log_2(K+d+N+2))\},\label{eq:runtime}
\end{align}
such that its sign-minimax risk is at most
\begin{equation}\label{eq:digitalrisk}
 (d/4)R_K^2+(\delta+\Gamma)^2.
\end{equation}
Here $W$ is additional temporary writable space, $L_{\rm prog}$ includes the schedule and coefficient table, and $T_{\rm bit}$ counts elementary operations on a deterministic sequential-tape implementation. Reads and arithmetic are padded to fixed phase lengths. A supplied clock with $Q\le C(N+T_{\rm bit}+1)$ phases is separately accounted for, giving total persistent cardinality at most $2^KQ$. Physical time is at most $N+T_{\rm bit}$ in the declared units. The terminal scored audit is after the output and is not a pre-decision digital input.

The same construction for Theorem~\ref{thm:joint} uses three additional bit/erasure labels, replaces $K$ by $K'$, and adds the attained bit risk $\varepsilon/4$ with the task weights. Choosing $p\ge\lceil\log_2((K+1)/R_K)\rceil$ puts numerical error at the matched geometric order. These are feasible upper accounts for workspace, program and time, not optimal lower bounds for those resources.
\end{theorem}
\begin{proof}
Append each received bit to a packed retained word and copy that word at idles. The scheduled coordinate and current position are determined by the padded phase, never by unretained observations. At the terminal decision, process coordinates successively. Start the accumulator at $1/2$, scan the retained word and the read-only schedule, and add $\widehat b_{j,l}(\zeta_{j,l}-1/2)$ for the selected coordinate. Dyadic arithmetic with a signed accumulator of $p+\lceil\log_2(K+2)\rceil+O(1)$ bits is exact for these approximate coefficients. Round the output to a $p$-bit dyadic and project to $[0,1]$.

Each coordinate's error from its exact conditional mean is at most $K2^{-p}/2+2^{-p}\le\Gamma$, and so is its weighted norm error. Conditional projection eliminates the cross term with $X-P(H)$. The perturbation and the inaccessible sign together have norm at most $\Gamma+\delta$, proving \eqref{eq:digitalrisk}. The output tape is write-only; later coordinates are computed from the retained word and calibration table, not by rereading previously emitted numbers.

The word uses exactly $K$ data bits. A coefficient and the signed accumulator, with loop indices, fit in \eqref{eq:workspace}. The table of at most $K$ coefficients, coordinate labels and finite loop parameters fits in \eqref{eq:program}. For each of $d$ coordinates perform at most $K+d$ fixed scans/additions, each with the displayed bit length; this deliberately loose sequential bound gives \eqref{eq:runtime}. Heads and loops follow predetermined padded phases, so their positions do not encode extra data. Counting those phases gives the bound on $Q$. A total-state constraint which includes the clock must use $2^KQ$, not just $2^K$. If the model instead supplies a deterministic external clock, its resource coordinate is still recorded. The three-state erasure extension is independent and gives the stated product account.
\end{proof}

\section{The ordinary filtering task}\label{sec:filters}
We verify the block side of Theorem~\ref{thm:transfer} for a structurally different raw experiment. The prediction loss is the usual squared error for a categorical hidden state, not an auxiliary centered probe designed to expose individual coordinates.

Fix $n=d+1\ge2$ and a finite action set. A known matrix $Q_t^a$ has a fixed primitive zero--one support $P$, with $P^m$ entrywise positive, and every nonzero entry of $Q_t^a$ is at least $s_*>0$. Thus every supported product of $m$ such matrices has all entries at least $s_*^m$. Prepare $Z_0$ uniformly on $\{0,\ldots,d\}$. After each transition, state $i$ emits the missing symbol $\bot$ with probability $1-q_i^a$ and otherwise a Gaussian vector of mean $\mu_i^a$ and covariance $\sigma^2I_d$. Require
\begin{equation}\label{eq:gaussparams}
 0<q_-\le q_i^a\le q_+<1,\quad \norm{\mu_i^a}_\infty\le R,
 \quad \sigma>0,
 \quad \abs{\det B_a}\ge b_*>0,
\end{equation}
where the $i$th row of $B_a$, $1\le i\le d$, is $(\mu_i^a-\mu_0^a)^{\mathsf T}/\sigma^2$. The determinant condition is for the acquired lower bound, not for the upper. The finite class parameters and primitive support are fixed; no assertion is uniform as these lower bounds vanish. Time-dependent known emissions can also be used when the same bounds hold.

With respect to Lebesgue measure and a unit atom at $\bot$ the emission density is
\[
 g_i^a(y)=q_i^a(2\pi\sigma^2)^{-d/2}
       e^{-\norm{y-\mu_i^a}^2/(2\sigma^2)},\qquad
 g_i^a(\bot)=1-q_i^a.
\]
For any fixed report-adapted legal exploration, the posterior is positive and has the raw Bayes update
\begin{equation}\label{eq:filterupdate}
 F_t(\pi,a,y)_i=
 \frac{(\pi Q_t^a)_i g_i^a(y)}{\sum_j(\pi Q_t^a)_j g_j^a(y)}.
\end{equation}
The posterior with its declared time/action interface determines all future tests. A terminal audit of $Z_T$, executed after the forecast, separates posterior coordinates, so this is a direct predictive realization. Every transition/report attempt, including a missing report, costs one call. Including preparation and the audit, raw acquisition is $T+2$.

Put $d_H(\pi,\rho)=\osc_i\log(\pi_i/\rho_i)$ and $D_H(\pi)=\osc_i\log\pi_i$. This is a separable metric on the open simplex with its usual Borel structure. The subscripts distinguish spread from the admissibility relation and controller cardinality.

\begin{lemma}[A mass-preserving finite cover]\label{lem:filtercover}
For some $\eta>0$ chosen from the raw class constants below, put $w(\pi)=e^{\eta D_H(\pi)}$. There are at most $M$ Borel code representatives with
\[
 d_H(\pi,Q_M\pi)\le C_{d,\eta}M^{-1/d}w(\pi),\quad
 D_H(Q_M\pi)\le D_H(\pi),\quad (Q_M\pi)_i\ge\pi_i/n.
\]
The uniform seed is a representative. The cover preserves the budget-independent relation
\begin{equation}\label{eq:filterrelation}
 \mathcal A(\pi)=\{\rho:D_H(\rho)\le D_H(\pi),\ \rho_i\ge\pi_i/n\text{ for all }i\}.
\end{equation}
\end{lemma}
\begin{proof}
Choose the least maximal coordinate $j$, and set $v_i=\log(\pi_j/\pi_i)\ge0$. For integer $K\ge1$ define
\[
 \widehat v_i=-\eta^{-1}\log\bigl(\lceil Ke^{-\eta v_i}\rceil/K\bigr),
 \qquad (Q\pi)_i=e^{-\widehat v_i}/\sum_h e^{-\widehat v_h}.
\]
At least one deficit is zero, and there are at most $nK^d$ representatives. The cells are Borel. From $u\le\lceil Ku\rceil/K\le u+K^{-1}$ one obtains
$0\le v_i-\widehat v_i\le e^{\eta v_i}/(\eta K)$.
The oscillation of these errors bounds the projective distance. Inwardness follows from $0\le\widehat v_i\le v_i$. Moreover $\sum_i e^{-v_i}\ge1$ and $\sum_i e^{-\widehat v_i}\le n$, proving the componentwise lower bound. Choose $K=\lfloor(M/n)^{1/d}\rfloor$ when $M\ge n$. For $M<n$ use only the uniform center and $D_H(\pi)\le e^{\eta D_H(\pi)}/(e\eta)$, increasing the fixed constant. The construction and relation are valid on the whole candidate simplex, not only on an observed trajectory.
\end{proof}

\begin{lemma}[Physical block contraction and robust drift]\label{lem:filterblock}
The update \eqref{eq:filterupdate} and relation \eqref{eq:filterrelation} satisfy \eqref{eq:pair}--\eqref{eq:within} with $L=1$, some $\kappa,a<1$ and finite $b,A,B$. Constants depend only on the fixed class parameters and are uniform in deterministic horizon and report-adapted actions. Individual updates need not contract strictly.
\end{lemma}
\begin{proof}
For a nonnegative matrix with no zero column, the derivative of $\log\sum_i e^{z_i}Q_{ij}$ is a convex average of the input direction. Its oscillation cannot increase. Integrating along a segment in log coordinates proves projective nonexpansiveness. Multiplication by a common positive likelihood vector cancels in the log-ratio difference, including at $\bot$.

Choose a fixed cube $G=[-b_1,b_1]^d$. On $G$, uniformly in the hidden state and action, $0<g_-\le g_i^a(y)\le g_+<\infty$. Given every past, $\PP(Y_t\in G\mid\mathcal H_{t-1})\ge p_0=g_-|G|>0$. If $E_j$ denotes the first $j$ reports of a block lying in $G$, then
\[
 \PP(E_j\mid\mathcal H_s)
 =\E[\one_{E_{j-1}}\PP(Y_{s+j}\in G\mid\mathcal H_{s+j-1})\mid\mathcal H_s]
 \ge p_0\PP(E_{j-1}\mid\mathcal H_s).
\]
Thus the good block $E_m$ has conditional mass at least $p=p_0^m$, without an independence assumption. Its unnormalized product matrix has entries in $[l,u]$, where $l=s_*^m g_-^m$ and $u=g_+^m$. Every column derivative distribution then dominates $(l/u)$ times the same input probability distribution. Removing that common mass gives projective gain $\chi=1-l/u<1$. Off the block use nonexpansiveness, so $\kappa=1-p(1-\chi^2)<1$.

For any adapted relation-preserving perturbed sequence, the good block gives componentwise
$z_j\ge n^{-1}(g_-/g_+)z_{j-1}Q_{s+j}^{a_{s+j}}$.
Its final coordinates are at least $c_*=n^{-m}(g_-/g_+)^ms_*^m$. Therefore $D_H(z_m)\le R_*:=\log(1/c_*)$. This proves why a cone-mass constraint, not spread inwardness alone, is needed through a sparse product.

For an arbitrary report let $V_a(y)=\osc_i\log g_i^a(y)$ and $C_s=\log(n/s_*)$. Then $D_H(F(z,a,y))\le D_H(z)+C_s+V_a(y)$, and admissible rounding does not increase spread. Gaussian log-likelihood ratios are affine functions of $y$, and the missing ratios are bounded. Hence for any fixed $\eta_0>0$,
$\mathcal M:=\sup_{x,a}\int e^{2\eta_0 V_a(y)}J_a(dy\mid x)<\infty$.
For example $V_a(y)\le v_0+v_1\norm{y}_1$, and
$\E e^{c\norm{Y}_1}\le2^d e^{cdR+dc^2\sigma^2/2}$ gives a uniform finite bound on the continuous component. Conditional iteration yields $\E e^{2\eta_0 Z}\le C_0$ for the block sum $Z=\sum_j(C_s+V_{a_{s+j}}(Y_{s+j}))$.

Choose $0<\eta\le\eta_0$ so that $(\eta/\eta_0)(C_0-1)\le p/2$. Convexity in the exponent gives
$e^{2\eta Z}-1\le(\eta/\eta_0)(e^{2\eta_0Z}-1)$.
Consequently $\E[e^{2\eta Z}\one_{E_m^c}\mid\mathcal H_s]\le1-p/2$. On $E_m$ the envelope is at most $e^{\eta R_*}$; off it, it is at most $w(z_0)e^{\eta Z}$. This proves \eqref{eq:drift} with $a=1-p/2$, $b=e^{2\eta R_*}$. The same growth estimate for every prefix gives \eqref{eq:within} with $A=C_0$, $B=0$. The proof applies to every adapted admissible perturbation under the true report law, not just to the chosen encoder.
\end{proof}

\begin{theorem}[Sharp state order for squared filtering loss]\label{thm:brier}
Under the preceding raw assumptions, fix $T\ge m$ and any report-adapted exploration independent of the predictor. Let
\[
 B_T^{\rm fil}=\E\norm{e_{Z_T}-\pi_T}_2^2.
\]
For prediction of the categorical state under loss $\norm{e_{Z_T}-A}^2_2$, with $A$ in the closed simplex,
\begin{equation}\label{eq:brier}
 cM^{-2/d}\le R_\cp(T,M)-B_T^{\rm fil}
 \le R_\on(T,M)-B_T^{\rm fil}\le CM^{-2/d}.
\end{equation}
The constants are uniform in $T$ and the specified exploration at fixed dimension, support, primitive length and parameter bounds. The readout state is a prediction label, not an unrestricted optimized controller. Controller and deterministic clock resources remain separate.
\end{theorem}
\begin{proof}
Differentiating the softmax along a log-coordinate segment gives derivative components $\pi_i(h_i-\sum_j\pi_jh_j)$. The sum of their absolute values, and hence their Euclidean norm, is at most $\osc(h)$. The posterior readout is therefore $1$-Lipschitz from $d_H$ to Euclidean norm. Lemmas~\ref{lem:filtercover} and \ref{lem:filterblock}, with the exact uniform seed, verify every upper hypothesis of Theorem~\ref{thm:transfer} and give $CM^{-2/d}$.

For the lower, retain only the actual event that the last $m$ reports lie in $G$. It has mass at least $p$. Conditional on the prefix and the first $m-1$ reports of this event, the predictive weights $r=\pi_{T-1}Q_T^{a_T}$ have coordinates at least $c_r=(g_-/g_+)^{m-1}s_*^m$. The last action is selected \emph{before} the final report. Its log posterior ratios are $B_{a_T}y$ plus a fixed vector, so the map from $y$ to the first $d$ posterior coordinates is injective. The softmax derivative in those coordinates is $\operatorname{diag}(\pi_1,\ldots,\pi_d)-vv^{\mathsf T}$, where $v=(\pi_1,\ldots,\pi_d)$. The matrix determinant lemma gives determinant $\prod_{i=0}^d\pi_i$. Thus the report-to-posterior Jacobian has absolute determinant at least $b_*c_p^{d+1}$ on $G$, where $c_p=c_r g_-/g_+$.

The last conditional report density is at most $g_+$. The resulting posterior submeasure has density at most $g_+/(b_*c_p^{d+1})$ in its first $d$ coordinates. Integrating over the conditioning prefix with its event indicator preserves this bound. The submeasure has its actual mass at least $p$; it is not normalized to mass one. An ambient Euclidean ball projects to a $d$-ball of the same radius, so its mass is at most $C_1r^d$. The union-of-balls lower in Theorem~\ref{thm:transfer} gives $cM^{-2/d}$. Projection for the categorical task identifies the common baseline and completes the proof.
\end{proof}

Report continuity is also explicit: mixtures of the same emission kernels have total variation at most $\frac12\norm{\pi Q-\rho Q}_1\le\frac12d_H(\pi,\rho)$. The terminal categorical audit obeys the same bound without $Q$. The total stopped-law comparison in Section~\ref{sec:composition} therefore applies. This risk theorem does not require a preassigned Fisher geometry; the exponent comes from a positive submeasure of the actual posterior law.

For a finite numerical interface, a bounded input box is essential. If all continuous reports before $T$ are clipped outside $[-U,U]^d$, then the failure probability is at most $2dT\exp(-(U-R)^2/(2\sigma^2))$. Choosing
\begin{equation}\label{eq:clipping}
 U=R+\sigma\sqrt{2\log(2dT/\zeta)}
\end{equation}
limits this defect to $\zeta\in(0,1)$. In particular $\zeta=M^{-2/d}$, with a harmless fixed reduction when $M=1$, requires growth of order $\sqrt{\log T+\log M}$ at fixed parameters. If intervals $[l_i,u_i]$ enclose unnormalized posterior log weights with width at most $h$, deficits $d_i=\max(0,\max_jl_j-u_i)$ satisfy $0\le d_i\le v_i$ and $v_i-d_i\le2h$. Applying the compander to these lower deficits preserves \eqref{eq:filterrelation} and gives local allowance $u=2h$. Proposition~\ref{prop:repair} charges clipping by bounded loss. These are enclosure and defect certificates; unlike Theorem~\ref{thm:digital}, they do not assert a complete optimal Gaussian bit-complexity region.

\section{Stopped interfaces and comparison of experiments}\label{sec:composition}
The lower and upper sides must refer to the same task under any experimental reduction. A complete-law defect is stronger than a predictive discrepancy and is added after root metric errors have been combined.

\subsection{Bounded stopped prediction}
For a visible stopping time $\tau\le mK_0$, define the stopped sigma field $\mathcal H_\tau$, a declared bounded task $Y^\tau$, and
\[
 P_\tau=\E[Y^\tau\mid\mathcal H_\tau],\qquad
 B_\tau=\E\norm{Y^\tau-P_\tau}^2.
\]
At $\tau=0$ these mean the actual preparation and its legal terminal task. Assume $P_\tau=f_\tau(S_\tau)$, with a common Lipschitz constant, and that every stopped exploration is the prefix of a specified legal continuation to $mK_0$ for which the block certificate holds. This continuation is used only for estimates; post-stop calls and information are not used by the actual predictor.

\begin{proposition}[A selection-sensitive stopped bound]\label{prop:stopped}
Under the upper hypotheses of Theorem~\ref{thm:transfer}, put $\rho=r_M+u$ and $\sigma=\lceil\tau/m\rceil$. There are explicit finite constants $c,C$, depending only on the block certificate, such that with
$V_0=d(S_0,\widehat S_0)^2+c\rho^2w(\widehat S_0)^2$,
\begin{equation}\label{eq:stoppedbound}
 \E d(S_\tau,\widehat S_\tau)^2
 \le C\{\E V_0+\rho^2\E\sigma\}.
\end{equation}
Consequently $e_0=O(\rho)$ and $\E\sigma\le L_0$ give
$R\le B_\tau+(C'\rho\sqrt{1+L_0}+v)^2$.
If the terminal decoder has a prediction label of size $M$ and a retained stopping clock of size $Q$, its checkpoint identity uses $q_{MQ}(\law(P_\tau))$ when this is the whole accessible tuple. A free declared interface instead requires conditioning on that interface.
\end{proposition}
\begin{proof}
Here are constants and the optional-sum argument, to distinguish the claim from substituting a random time into a deterministic estimate. Put
\[
 \bar\kappa=(1+\kappa)/2,\quad C_Y=\bar\kappa/(\bar\kappa-\kappa),
 \quad\gamma=\max\{\bar\kappa,(1+a)/2\},
\]
\[
 c=\max\{1,C_Y\Lambda_m^2A/(\gamma-a)\},\qquad
 c_0=C_Y\Lambda_m^2(B+1)+cb.
\]
The conditional shadow bound, $r_M\sqrt{x}+u\le\rho\sqrt{x+1}$ and
$(\sqrt\kappa d+z)^2\le\bar\kappa d^2+C_Yz^2$ give
\[
 \E[V_{k+1}\mid\mathcal H_{km}]\le\gamma V_k+c_0\rho^2,
 \quad V_k=d(S_{km},\widehat S_{km})^2+c\rho^2w(\widehat S_{km})^2.
\]
Multiply the rearranged inequality by the predictable indicator $\one_{\{k<\sigma\}}$ and sum the finitely many terms. This gives
$(1-\gamma)\E\sum_{k<\sigma}V_k\le\E V_0+c_0\rho^2\E\sigma$.
For a raw stopping time, the within-block telescoping bound and the inequality that a maximum of nonnegative terms is at most their sum give
\[
 \E[\max_{1\le j\le m}d(S_{km+j},\widehat S_{km+j})^2\mid\mathcal H_{km}]
 \le C_1V_k+C_2\rho^2,
\]
where $C_1=2m(L^{2m}+\Lambda_m^2A/c)$ and $C_2=2m\Lambda_m^2(B+1)$. Sum only over entered blocks; their indicator is known at each block start. At $\tau=0$ use $V_0$. The resulting bound is
\[
 (1+C_1/(1-\gamma))\E V_0+
 (C_1c_0/(1-\gamma)+C_2)\rho^2\E\sigma,
\]
which proves \eqref{eq:stoppedbound}. Stopped conditional projection proves the risk statement. The information-cut identity follows by counting every available terminal state before applying Lemma~\ref{lem:cut}; it is not an unqualified $M$-center formula when the clock carries information.
\end{proof}

\subsection{Measurable report coupling}
For probability kernels $P_z,Q_z$ on a standard Borel report space, a jointly measurable maximal coupling exists. One construction uses nested finite partitions generated by a countable separating algebra. Relative to $R_z=(P_z+Q_z)/2$, their atomwise density ratios are bounded measurable functions. Martingale convergence gives a jointly measurable clipped limsup version $p=dP_z/dR_z$ for each $z$; put $q=2-p$. The common kernel $\mu_z=\min(p,q)R_z$ has mass $c_z$. The coupling
\[
 \operatorname{diag}_*\mu_z+
 \frac{(P_z-\mu_z)\otimes(Q_z-\mu_z)}{1-c_z}\quad(c_z<1)
\]
and the diagonal law for $c_z=1$ are measurable kernels. The residual measures are mutually singular, so their product has zero diagonal mass. The disagreement probability is $1-c_z=\TV(P_z,Q_z)$, with $\TV=\sup_A|P(A)-Q(A)|$.

Suppose report kernels satisfy $\TV(J(\cdot\mid x,a),J(\cdot\mid z,a))\le L_Jd(x,z)$, and task kernels have constant $L_Y$. A preparation mismatch is either absent or explicitly at most $\varepsilon_0$. Successive maximal couplings, driven until first disagreement by the same declared controller and stopping rule, give complete stopped report-and-task defect at most
\begin{equation}\label{eq:reporttv}
 \min\{1,\varepsilon_0+L_J\E\sum_{t=1}^\tau d(S_{t-1},\widehat S_{t-1})
                         +L_Y\E d(S_\tau,\widehat S_\tau)\}.
\end{equation}
Indeed before disagreement both transcripts, actions and stopping decisions coincide. Continue a physical-law shadow after disagreement and sum first-disagreement probabilities, dropping the agreement indicator. This preserves the unconditional shadow estimates. A stopped sum is not replaced by $\E\tau$ times an unconditional deterministic-time bound. A finite complete-law bound is not infinite-path or large-deviation equivalence.

\subsection{Resource-certified morphisms}
For clarity we state all components of the comparison used here. A source experiment $\mathfrak E$ implements a target $\mathfrak F$ on a common parameter set by parameter-independent transducer initialization, state-update, source-action and report-transformation kernels. These kernels induce a strategy lift from the declared target class to legal source strategies. Each virtual completion time $\sigma_i$ is a source stopping time bounded by the given physical budget; for every permitted target stop $\tau$, $\sigma_\tau$ is a legal source stop and no later source information enters its decision. Legal actions and the available interface agree after each transformed prefix. Virtual costs, preparation and failures are included.

There is one Borel common scored-outcome map and one decision space, with loss in $[0,L_*]$. Uniformly over the common parameters, declared target policies and permitted bounded stops, the joint law of virtual history, permitted public policy coins and scored outcome differs from its target law by at most $\varepsilon$ in total variation. A transducer relying on a real hidden history or unspecified phase is not a finite implementation. Let its state count be $R$, target prediction count $M$, controller count $C_{\rm ctr}$ and phase/clock count $Q$. Workspace, read-only program, precision, raw calls and physical time are separate coordinates. The companion~\cite{companion} gives the fully typed standard Borel transcript construction; these requirements are precisely the ones needed in the following implication.

\begin{theorem}[Common-risk composition and reverse transfer]\label{thm:morphism}
Such a certificate implements a target risk $r_\lambda$ by a serial source tuple of at most $MR C_{\rm ctr}Q$ states and risk at most $r_\lambda+L_*\varepsilon$. Defects of two composable certificates add, capped at one, provided the lifted policies, stopping rules and common-task maps lie in the next certified class. Serial state overheads multiply; raw costs sum over executed subroutines, and nested per-call upper bounds may multiply.

If every source $K$-state procedure can conversely be implemented in the target with at most $\Psi(K)$ states and complete defect $\varepsilon_-$, then for the same task and risk criterion
\begin{equation}\label{eq:morphismlower}
 R_E(K)\ge R_F(\Psi(K))-L_*\varepsilon_-.
\end{equation}
In particular every acquired continuation lower of Lemma~\ref{lem:cut} transfers with that inflated budget and the target baseline left unchanged. A block upper transfers as
\begin{equation}\label{eq:morphismupper}
 R_E(MRC_{\rm ctr}Q)
 \le B_T+(L_fE_{M,u}+v)^2+L_*\varepsilon.
\end{equation}
The same assertions hold for minimax risks when both certificates are uniform over the same parameter set.
\end{theorem}
\begin{proof}
Run the physical subroutine and retain the tuple of predictor, simulator, controller and clock states. The action and stopping conditions make the procedure executable; its cardinality is the displayed product. Applying the common decision kernel preserves total-variation closeness. The bounded-loss inequality for our convention gives a difference at most $L_*\varepsilon$.

For composition insert the intermediate stopped joint law and use the triangle inequality and contraction of total variation under the downstream parameter-independent transformation. Class closure is needed here for adaptive lifted policies. Retaining both serial tuples proves the storage upper. All actual calls, waiting times and failures are summed; shared program or workspace can sometimes be reused, so additive serial descriptions are only upper accounts.

For the reverse statement apply the upper implication to every source procedure: $R_F(\Psi(K))\le r_E+L_*\varepsilon_-$. Infimize over source procedures. Uniformity permits the same argument after taking the parameter supremum. No separately recomputed source and target Bayes baselines are subtracted. This proves both \eqref{eq:morphismlower} and \eqref{eq:morphismupper} in their stated directions.
\end{proof}

The exact-bit and erasure pair in Theorem~\ref{thm:joint} verifies this definition with a prescribed completion clock, parameter-independent bit transformation and a common terminal audit; it is not just a marginal report comparison. The Gaussian report comparison verifies \eqref{eq:reporttv} for the same declared controller. These are distinct uses: the latter does not manufacture a reverse certificate for all Gaussian source algorithms. Likewise a product $RC_{\rm ctr}Q$ is not necessary merely because one implementation stores it. Only an executable challenge separating those registers, or an appropriate reverse theorem, provides that lower bound.

\section{Antecedents, distinctions and quantitative scope}\label{sec:literature}
\subsection{What the transfer theorem adds}
Conditional projection, small-ball quantization and finite-center approximation are classical; see Graf and Luschgy~\cite{graf}. Predictive representations and future equivalence are developed by Littman, Sutton and Singh~\cite{psr} and Shalizi and Crutchfield~\cite{shalizi}; measurable kernel sufficiency has a categorical treatment in Fritz~\cite{fritz}. Blackwell~\cite{blackwell}, Le Cam~\cite{lecam} and Torgersen~\cite{torgersen} provide the statistical-comparison background. Proposition~\ref{prop:quotient} makes explicit a measurable-image certificate and common report versions; it does not originate predictive sufficiency. The classical total-variation inequality is not the new part of Theorem~\ref{thm:morphism}.

Average contraction of random maps is also classical~\cite{diaconis}. The constructive part of Theorem~\ref{thm:transfer} uses a stronger, specifically physical-law certificate: a cover-preserved relation must control every adapted candidate perturbation across a block. The positive-cone lower bound in Lemma~\ref{lem:filtercover} shows how such a relation can be manufactured for primitive sparse products. Its role is to preserve accumulated positivity while rounding after each raw call. It is sufficient, not an intrinsic necessary condition. Proposition~\ref{prop:repair} allows flagged implementation failures, but does not show that arbitrary local covers or ordinary inwardness imply robust block drift.

The opposite direction now uses actual distributions of continuation-response functions under a joint-law domination certificate. Unlike exact tag separation, it allows a continuum of noisy response functions and charges distortion in the same terminal squared task. Unlike a last-marginal small-ball estimate, it can distinguish a large information cut from an atomic terminal answer. The proof is a Hilbert-space quantization argument; the random-access example explicitly uses an established entropy mechanism~\cite{nayak}. The companion's earlier exact-tag and soft-information arguments remain complementary, rather than being renamed as new theorems.

The optimized acquisition result uses a different constructive mechanism: exact refinement of acquired prefixes, with no rounding of an unretained real posterior. Quantization on Moran sets and measures has an extensive independent theory; for example Zhu~\cite{zhu-moran} proves comparability of optimal Voronoi-cell errors under doubling and separation assumptions. Lemma~\ref{lem:greedy} is a self-contained nonasymptotic product estimate. Its use in Theorems~\ref{thm:adaptive} and \ref{thm:joint} is to identify one scale which also controls every adaptive allocation of paid digit queries. We do not claim a general new quantization theorem for all Moran measures, nonproduct sources or arbitrary adaptive partitions.

\subsection{Operationally close work}
Recursive nonlinear-filter quantization predates this work. Pag\`es and Pham~\cite{pages} and Sellami~\cite{sellami} transport quantization error through filtering recursions. Their approximation constructions do not by themselves identify a fixed cardinality for every retained predictive label or a matching actual-posterior converse. Here Lemmas~\ref{lem:filtercover} and \ref{lem:filterblock} provide that count, and Theorem~\ref{thm:brier} gives the ordinary filtering excess. No novelty is claimed for the Bayes formula, softmax derivative or projective nonexpansiveness separately.

Theorem~9 of Subramanian, Sinha, Seraj and Mahajan~\cite{ais} bounds approximate dynamic-program values from an approximate information-state generator and reward/transition discrepancies. Its policy-loss conclusion concerns control optimization. Theorem~16 of Kara and Y\"uksel~\cite{kara} derives finite-window control bounds under their stability, finite-alphabet and discount hypotheses. Our same-controller physical-law estimate does not imply either optimal-control conclusion. Conversely their conditions do not supply the actual acquired-law lower bound or the hard predictive-state alphabet of \eqref{eq:brier}.

Reachable-belief covering is not a new complexity viewpoint. Zhang, Littman and Chen~\cite{zhang2012} connect covering numbers to POMDP planning and no-reset learning. Theorem~1 of Zhang, Hsu and Lee~\cite{zhang2014} gives a heuristic-guided planning-time upper of order $h\,C(\delta/2)^2$, with depth and resolution chosen from the discount and target value error. Its cover concerns a heuristically reachable search space, not the probability mass actually acquired under a fixed experiment. Grover and Dimitrakakis~\cite{grover} provide adaptive belief discretization and explicit planner-memory/value-error bounds. They therefore cannot be dismissed as not addressing memory; the distinction is planning storage and value approximation versus per-cut predictive alphabet, acquisition uncertainty and a matching geometric lower.

Recent work likewise uses belief geometry operationally. Zhu and Lu~\cite{zhulu} develop offline POMDP evaluation bounds through belief-space coverage and policy stability. This is a learned off-policy evaluation problem, whereas \eqref{eq:central} fixes known kernels and \eqref{eq:adaptivelaw} optimizes a specified progressive sensor. Anjarlekar, Etesami and Srikant~\cite{anjarlekar} use finite-history superstates and policy-based learning with approximation controlled by history length. We do not obtain their learning guarantees from our known-kernel digital implementation. These comparisons are to their published conference versions, not predictions of future results.

Finite-state controller synthesis also has an independent objective. Andriushchenko, \v{C}e\v{s}ka, Junges and Katoen~\cite{inductive} synthesize controllers for POMDP specifications. Hud\'ak and collaborators~\cite{hudak} extract finite-state controllers from learned recurrent policies and evaluate them, including hidden-model robustness. These are substantial finite-controller results, not instances of a terminal codebook. Our converse fixes a common statistical task and the information crossing a cut; it does not establish controller synthesis completeness, robustness for unknown dynamics, or superiority in computational time.

Zero-delay coding and nonanticipative rate-distortion theory address related causal constraints. Wood, Linder and Y\"uksel~\cite{wood}, Ghomi, Linder and Y\"uksel~\cite{ghomi}, and Charalambous, Stavrou and Ahmed~\cite{nonanticipative} study structural or information-constrained formulations. A mutual-information budget is not automatically a hard retained alphabet; nor does our finite-horizon state-order result prove existence of optimal stationary average-cost codes.

\subsection{Scope of the resulting laws}
The new obstruction and the sufficient construction deliberately have different assumptions. Joint-law domination may have very small mass or fail entirely. Passing all the available cut tests is not proved sufficient for a compatible causal implementation. The profile \eqref{eq:adaptivelaw} is optimized over every adaptive allocation in its raw sensing class, but not over arbitrary interventions, unknown-kernel learning or unbounded stopping. It includes only the stated deterministic phase convention; total persistent state includes any separately stored clock.

For the Gaussian realization, constants deteriorate with $1-\sqrt\kappa$, $1-a$, primitive length, positivity floors, acquired mass and the inverse determinant in \eqref{eq:gaussparams}. For the multiscale realization they depend on $a,d$, not on a nonexistent assumed limiting dimension. Theorems~\ref{thm:digital} and \ref{thm:morphism} charge workspace, program, simulator and time, but do not prove the full optimal feasible region for those coordinates. A lower bound on an output grid is not a lower bound on each internal arithmetic operation. A calibration or deficiency \emph{allowance} alone is not a floor; \eqref{eq:jointlaw} attains such terms only in its explicit inaccessible-sign and erasure family.

Thus the central connection is mathematical rather than terminological: the prepared kernels determine future tests, their measurable predictive realization and actual acquired laws; compatible finite recursions construct the upper; dominated continuation cuts and acquisition variance supply converses; and typed causal morphisms transfer both at their declared resources. A necessary-and-sufficient classification of all recursive acquired geometries remains beyond the theorems proved here.

\appendix
\section{Retained singular dynamics and interface conventions}\label{sec:appendix}
\subsection{The nonreset singular realization}
The technical companion~\cite{companion}, Section~5, supplies a second direct verification of the block hypotheses, distinct from Gaussian filtering. We record its raw data and the exact point of attachment so that its status is not confused with the adaptive progressive-sensor class of Section~\ref{sec:adaptive}.

Take ratios $r_l\in[1/5,1/3]$, $\ell_l=\prod_{j\le l}r_j$ and a fair binary Cantor coordinate $x=\sum_l(1-r_l)\ell_{l-1}\zeta_l$. Complete its domain by adjoining every finite-prefix fair-tail mean. A prepared mode $I\in\{0,1\}$ is reported, followed by $B$ paid initial prefix reads. At a later opportunity an independent active indicator has any deterministic probability. Inactivity is an exact copy but costs a call. On activity, mode zero either prepends two fair reported bits, deletes the first bit within mode zero, or changes to mode one while deleting; the probabilities are $(1-\alpha_t)(1-\gamma_0)$, $(1-\alpha_t)\gamma_0$, and $\alpha_t$, with $\gamma_0=1/4096$. Mode one prepends six reported fair bits and moves to zero with probability $1/2$, or copies within mode one with probability $1/2$. Deletion does not reveal the deleted bit; both mode changes transport a nonconstant old tail.

Within each mode give the completed coordinate Euclidean distance weighted by $v_0=1800$, $v_1=1$; between modes use distance $2\sqrt{1800}$. The relation is the full same-interface candidate domain and $w=1$. On the completed domain deletion has gain at most $30$ and prepending $s$ fixed bits has gain at most $6\,3^{-s}$. The matrix of physical expected squared gains is
\[
 C_t=\begin{pmatrix}(1-\alpha_t)85/128&900\alpha_t\\2/59049&1/2\end{pmatrix},
 \qquad C_t\binom{1800}{1}\le\frac{17}{20}\binom{1800}{1}.
\]
Cross-mode pairs have squared gain at most $1/4$. Thus the active block certificate is $m=1$, $\kappa=17/20$, and constant envelope. All finite-prefix centers through depth $h$ in both modes form a global Borel cover of cardinality $2^{h+2}-2$ and radius $\sqrt{1800}\ell_h$. The gain estimates follow by comparing the first distinguishing level: its original separation lies between $\ell_{k-1}/6$ and $\ell_{k-1}$, while deletion or prepending shifts this level. These facts verify the hypotheses independently of any claimed memory bound.

The companion proves the resulting exact finite-acquisition law, including the physical-oracle lower and the mode-dependent affine terminal Bernoulli task. If $H_T$ is the actual acquired depth, then the actual score law is a \emph{finite atomic} mixture of the fair-prefix means. Conditional tail variance is comparable to $\E\ell_{H_T}^2$, and a physical Cantor small-ball witness gives the memory term $\ell_{\lfloor\log_2M\rfloor}^2$. Their maximum and sum are both comparable to
\[
 \E\ell_{\min(H_T,\lfloor\log_2M\rfloor)}^2.
\]
The matched sign-minimax law adds $\delta^2$ and charges $B+T+2$ raw calls. The physical law is used solely as an oracle lower witness, never substituted for the actual atomic acquired law. Its proof and all additional assumptions are retained in full in the companion. The progressive acquisition theorem here is a different result: it permits adaptive allocation across coordinates and even smooth factors, but does not replace this nonreset dynamical theorem.

For completeness, exact independent holds preserve the active error radius without requiring a uniform positive activation probability. If active updates satisfy $e_+\le\sqrt\kappa e+\Delta$ and holds copy exactly with independent probability $1-\vartheta_t$, then
\[
 e_t^2\le(1-\vartheta_t)e_{t-1}^2+
       \vartheta_t(\sqrt\kappa e_{t-1}+\Delta)^2.
\]
Both terms preserve $e\le\max\{e_0,\Delta/(1-\sqrt\kappa)\}$. The envelope argument is identical with indicator-form expectations. No conditioning on a null active event is needed when $\vartheta_t=0$. State-dependent missingness in Section~\ref{sec:filters} is not such an independent hold.

\subsection{Finite interfaces and the location of cuts}
A publicly chosen terminal index after an encoding cut means that each retained label selects a vector, or a response function, of possible readouts. It does not mean that the encoder saw the selected index before that cut. This is why the continuation function in Example~\ref{ex:query} has dimension $d$ while its terminal answer has only two values. When a clock or an interface tag can be queried after the cut, its information is part of the retained tuple unless the theorem explicitly conditions on that public interface.

An exact continuation challenge can force disjoint state supports. For a tag $J$ with probabilities $p_j$, unreissued after the cut, exact recovery forces at least one disjoint label set of size $m_j$ per tag. If the score's conditional law is $\nu_j$, the precise distortion is
\[
 \inf_{m_j\ge1,\ \sum_jm_j\le K}\sum_jp_jq_{m_j}(\nu_j).
\]
The proof fixes independent coins, observes that one label cannot decode two distinct tags exactly, and applies conditional nearest-center quantization within each support. Disjoint codebooks give the reverse inequality. This inherited tag mechanism is different from the noisy continuation-function bound of Lemma~\ref{lem:cut}; neither licenses an unchallenged simulator-state lower bound.

On the analytic side all report versions used by wrong candidates must be jointly Borel on the declared true/candidate domain. In the Gaussian case emission densities are positive and the update is everywhere defined on the open simplex. In the continuous-instrument proposition only the positive-density domain is normalized. The root allowance $u$ in Theorem~\ref{thm:transfer} is valid only for routines still satisfying the relation, not for arbitrary floating-point perturbations of the same size. A calibrated-input perturbation can be included in $u$ only after proving that same stable-metric and relation certificate.

These conventions also delimit the retained companion statements. Finite-dimensional identities there do not imply unbounded-operator closure; finite stopped total variation does not imply a large-deviation principle; and same-controller coupling does not imply unrestricted optimal control. The independent realization articles on ordered measurements, finite physical actions and streaming space are not premises of the foundation theorem.

\begin{thebibliography}{99}
\bibitem{anjarlekar} A. Anjarlekar, S. R. Etesami and R. Srikant, \emph{Scalable policy-based RL algorithms for POMDPs}, Advances in Neural Information Processing Systems \textbf{38} (2025).
\bibitem{inductive} R. Andriushchenko, M. \v{C}e\v{s}ka, S. Junges and J.-P. Katoen, \emph{Inductive synthesis of finite-state controllers for POMDPs}, Proc. UAI, Proceedings of Machine Learning Research \textbf{180} (2022), 85--95.
\bibitem{blackwell} D. Blackwell, \emph{Equivalent comparisons of experiments}, Ann. Math. Statist. \textbf{24} (1953), 265--272.
\bibitem{nonanticipative} C. D. Charalambous, P. A. Stavrou and N. U. Ahmed, \emph{Nonanticipative rate distortion function and relations to filtering theory}, IEEE Trans. Automat. Control \textbf{59} (2014), 937--952.
\bibitem{diaconis} P. Diaconis and D. Freedman, \emph{Iterated random functions}, SIAM Rev. \textbf{41} (1999), 45--76.
\bibitem{fritz} T. Fritz, \emph{A synthetic approach to Markov kernels, conditional independence and theorems on sufficient statistics}, Adv. Math. \textbf{370} (2020), 107239.
\bibitem{ghomi} M. Ghomi, T. Linder and S. Y\"uksel, \emph{Zero-delay lossy coding of linear vector Markov sources: optimality of stationary codes and near optimality of finite memory codes}, IEEE Trans. Inform. Theory \textbf{68} (2022), 3474--3488.
\bibitem{graf} S. Graf and H. Luschgy, \emph{Foundations of Quantization for Probability Distributions}, Lecture Notes in Mathematics 1730, Springer, 2000.
\bibitem{grover} D. Grover and C. Dimitrakakis, \emph{Adaptive belief discretization for POMDP planning}, arXiv:2104.07276, 2021.
\bibitem{hudak} D. Hud\'ak, M. F. L. Galesloot, M. Tappler, M. Kure\v{c}ka, N. Jansen and M. \v{C}e\v{s}ka, \emph{Finite-state controllers for (hidden-model) POMDPs using deep reinforcement learning}, Proc. AAMAS (2026), 1201--1210.
\bibitem{kara} A. D. Kara and S. Y\"uksel, \emph{Near optimality of finite memory feedback policies in partially observed Markov decision processes}, J. Mach. Learn. Res. \textbf{23} (2022), no.~11, 1--46.
\bibitem{lecam} L. Le Cam, \emph{Asymptotic Methods in Statistical Decision Theory}, Springer, New York, 1986.
\bibitem{psr} M. L. Littman, R. S. Sutton and S. Singh, \emph{Predictive representations of state}, Advances in Neural Information Processing Systems \textbf{14} (2001), 1555--1561.
\bibitem{nayak} A. Nayak, \emph{Optimal lower bounds for quantum automata and random access codes}, Proc. 40th IEEE Symposium on Foundations of Computer Science (1999), 369--376.
\bibitem{pages} G. Pag\`es and H. Pham, \emph{Optimal quantization methods for nonlinear filtering with discrete-time observations}, Bernoulli \textbf{11} (2005), 893--932.
\bibitem{companion} Q. Qi, \emph{General Theta Foundations I: Acquired Geometry and Causal Resource Transfer}, accompanying technical article, 2026, Sections~1--7.
\bibitem{sellami} A. Sellami, \emph{Quantization based filtering method using first order approximation}, SIAM J. Numer. Anal. \textbf{47} (2010), 4711--4734.
\bibitem{shalizi} C. R. Shalizi and J. P. Crutchfield, \emph{Computational mechanics: pattern and prediction, structure and simplicity}, J. Stat. Phys. \textbf{104} (2001), 817--879.
\bibitem{ais} J. Subramanian, A. Sinha, R. Seraj and A. Mahajan, \emph{Approximate information state for approximate planning and reinforcement learning in partially observed systems}, J. Mach. Learn. Res. \textbf{23} (2022), no.~12, 1--83.
\bibitem{torgersen} E. Torgersen, \emph{Comparison of Statistical Experiments}, Cambridge University Press, 1991.
\bibitem{wood} R. G. Wood, T. Linder and S. Y\"uksel, \emph{Optimal zero delay coding of Markov sources: stationary and finite memory codes}, IEEE Trans. Inform. Theory \textbf{63} (2017), 5968--5980.
\bibitem{zhang2012} Z. Zhang, M. Littman and X. Chen, \emph{Covering number as a complexity measure for POMDP planning and learning}, Proc. AAAI \textbf{26} (2012), 1853--1859.
\bibitem{zhang2014} Z. Zhang, D. Hsu and W. S. Lee, \emph{Covering number for efficient heuristic-based POMDP planning}, Proc. ICML, Proceedings of Machine Learning Research \textbf{32} (2014), 28--36.
\bibitem{zhu-moran} S. Zhu, \emph{On the asymptotic quantization error for the doubling measures on Moran sets}, arXiv:1908.00202, 2019.
\bibitem{zhulu} Y. Zhu and Y. Lu, \emph{A covering framework for offline POMDPs learning using belief space metric}, Proc. AISTATS, Proceedings of Machine Learning Research \textbf{300} (2026), 1423--1431.
\end{thebibliography}

\end{document}
